% mode: ini
% Aimed at the linked hash of section 1117 and 1118: words of one hash value, chained through hyph_link above hyph_prime, and hyph_next counting up; a repeat found along a chain; the same letters in another language, which is
% a different string; and the lookup of section 1107 walking the chain, finding, and missing at the end of it.  a-ai a-bg a-ce a-dc a-ea b-ae all hash to the same slot for language 0.
\input prelude
\lefthyphenmin=1 \righthyphenmin=1 \defaulthyphenchar=`\- \hyphenpenalty=-10000 \hsize=1pt \tolerance=10000 \parfillskip=0pt \parindent=0pt
\tenrm \hyphenchar\tenrm=`\-
% Twice each: in a narrow measure, so that each hyphen breaks a line, and in a wide one, where the discretionaries stay whole and show pre_break, post_break and "replacing".
\def\try#1{{\hsize=1pt \setbox0=\vbox{\hskip0pt #1\par}\showbox0 }{\hsize=1000pt \hyphenpenalty=50 \linepenalty=10 \tolerance=-1 \setbox0=\vbox{\hskip0pt #1\par}\showbox0 }}
\hyphenation{a-ai a-bg a-ce a-dc a-ea b-ae}
\try{aai}
\try{abg}
\try{ace}
\try{adc}
\try{aea}
\try{bae}
% A word that hashes there but is not in: the chain ends.  aai with a letter changed gives another slot; this one has the same length.
\try{aag}
% Repeats along the chain, with a new position: first and last.
\hyphenation{aa-i}
\try{aai}
\hyphenation{ba-e}
\try{bae}
% One more on the end of the chain.
\hyphenation{ad-c ae-a}
\try{adc}
\try{aea}
% Words of other lengths in the slot of nothing else.
\hyphenation{aa-ii aa-iii aa-iiii}
\try{aaii}
\try{aaiii}
\try{aaiiii}
% Patterns are not used for a word with an exception.
\language=1
\hyphenation{a-ai}
\try{aai}
\try{abg}
\language=0
\try{aai}
\end
